\section{Diskussion}
\label{sec:discussion}

\subsection{Warum HAWQ-V2 unter PoT scheitert}
\label{sec:disc-hawq}

Die Taylor-Motivation hinter dem HAWQ-V2-Produkt (Gleichung~\eqref{eq:taylor}) setzt voraus, dass die Störungsgröße $\|\Delta W_l\|^2$ über die Schichten hinweg vergleichbar skaliert, sodass die Krümmung $\mathrm{Tr}(H_l)$ die eigentliche Rangfolge dominiert. Unter uniformer Quantisierung ist diese Annahme plausibel. Die Störungsgröße haengt dort im Wesentlichen von der gewählten Bit-Breite und dem Wertebereich einer Schicht ab, und beide Größen unterscheiden sich zwischen den Schichten eines Netzes oft nur moderat.
\par\medskip

Unter PoT bricht diese Annahme zusammen. Wie in \S\ref{sec:fund-quant} erläutert, hängt die PoT-Störungsgröße primaer davon ab, wie gut die spezifischen Gewichtswerte einer Schicht zufällig mit dem logarithmischen PoT-Gitter übereinstimmen, was eine weitgehend von der Kruemmung dieser Schicht unabhängige Eigenschaft ist. \texttt{conv1} hat in den meisten Kombinationen die höchste Sensitivität (höchster $S_{\mathrm{raw}}$-Rang, Tabelle~\ref{tab:rank}), aber eine der kleinsten Störungsgrößen, sodass die beiden Faktoren des HAWQ-V2-Produkts hier entgegengesetztes Verhalten zeigen. Damit korreliert $S_{\mathrm{pert}}$ in 5 von 6 Kombinationen negativ mit dem gemessenen Schaden,  entgegengesetzt zu $S_{\mathrm{raw}}$s positiver Korrelation in denselben 5 Kombinationen (\S\ref{sec:phase4}). Das Produkt $S_{\mathrm{hawq}}$ wird dadurch systematisch nach unten gezogen. Statt dass sich Kruemmungs- und Störungssignal wie in der Taylor-Motivation vorgesehen sinnvoll ergänzen, heben sie sich teilweise auf.
\par\medskip

Es ist kein allgemeines Scheitern Krümmungs-basierter Sensitivität, sondern ein spezifisches Scheitern der Produktbildung genau unter PoT. Das Produkt funktioniert nur dann als sinnvoller Sensitivitäts-Score, wenn $\|\Delta W\|^2$ zumindest unkorreliert mit der Krümmungs-basierten Sensitivität einer Schicht ist. PoT verletzt dies strukturell, da die Gitter-Ausrichtung eine Eigenschaft der Gewichtsverteilung und unabhängig von der Krümmung ist, und für \texttt{conv1} sind beide direkt gegenläufig. Dies erklärt, warum eine etablierte, breit genutzte Sensitivitätsformulierung spezifisch unter diesem Quantisierungsschema in die Irre fuehrt.

\subsection{Konfigurations-Sensitivität als methodisches Problem}
\label{sec:disc-config}

Dass zwischen den historischen und den kanonischen Werten derselben nominellen Größe zwei Groeßenordnungen Abweichung liegen, ohne dass eine einzelne Konfigurationsentscheidung im gesweepten Parameterraum sie erklärt, zeigt, wie fragil ein einzelner, nicht dokumentierter Trace-Wert als Forschungsartefakt ist. Publizierte Krümmungs-basierte Sensitivitätsanalysen, die ihre Konfiguration (Loss-Reduktion, BN-Fusionsstatus, Proben-Budget, Datensplit, Zufalls-Seeds) nicht explizit dokumentieren, sind mit hoher Wahrscheinlichkeit schwer bis gar nicht reproduzierbar. Dies liegt nicht daran, dass die zugrundeliegende Methode fehlerhaft wäre, sondern daran, dass der Wertebereich, den unterschiedliche plausible Konfigurationsentscheidungen erzeugen können, in den eigenen Messungen mehrere Größenordnungen umfasst (\texttt{drift\_diagnosis.csv}, \S\ref{sec:phase3}).
\par\medskip

Daraus lässt sich eine konkrete Empfehlung für zukünftige Arbeiten in diesem Bereich ableiten. Eine kanonische, explizit versionierte Konfigurationsdatei (wie die vorliegende \texttt{trace\_config.json}) sollte Standardpraxis sein, ebenso wie das Protokollieren, welche historischen Werte mit welcher Konfiguration reproduziert werden konnten und welche nicht, statt nur den zuletzt gemessenen Wert ohne diese Historie zu berichten.

\subsection{Gewichts- vs.\ Aktivierungsregime}
\label{sec:disc-regime}

Die Dekomposition aus \S\ref{sec:phase1-results} zeigt, dass unter PTQ in allen drei getesteten Architekturen der Gewichtsanteil den Genauigkeitsverlust auf CIFAR10 dominiert. Da die Hessian-Analyse ausschließlich Gewichtsschaden erfasst, deckt sie damit, zumindest fuer PTQ auf CIFAR10, den tatsächlich dominierenden Verlustmechanismus ab. Unter QAT kehrt sich das Bild um. Dort dominiert der Aktivierungsanteil, und eine Übertragung der Hessian-basierten Gewichts-Sensitivitätsanalyse auf diesen Fall wäre nicht gerechtfertigt. Diese Beobachtung, konsistent mit der grundsaetzlich unterschiedlichen Natur von Gewichts- und Aktivierungsschaden aus \S\ref{sec:fund-activation}, ist der Kern der Antwort auf die dritte Forschungsfrage. Die Hessian-basierte Sensitivitaetsanalyse ist nicht universell auf jedes Quantisierungsregime anwendbar, sondern regimeabhängig, und ihre Anwendbarkeit sollte vor jedem Einsatz explizit geprüft werden, etwa durch eine Dekomposition wie in \S\ref{sec:phase1-results}, statt vorausgesetzt zu werden.

\subsection{PTQ vs.\ QAT Krümmungsstruktur}
\label{sec:disc-ptqqat}

Die Ergebnisse aus \S\ref{sec:phase9} beantworten die erste Forschungsfrage mit einer auf den ersten Blick überraschenden Kombination. Einerseits verändert QAT die relative schichtweise Krümmungsstruktur eines Netzes kaum, die Spearman-Korrelation zwischen den PTQ- und QAT-Trace-Profilen liegt bei $\rho=0.95$--$1.00$ über alle sechs Kombinationen. Andererseits bricht die Beziehung zwischen dieser Krümmung und dem tatsächlich beobachteten Schaden unter QAT fast vollständig zusammen (0 von 18 loss-basierten Korrelationen signifikant, gegenüber 5 von 18 unter PTQ, \S\ref{sec:phase4}). Auf den ersten Blick widerspricht sich das, dieselbe Krümmungslandschaft sollte doch eine ähnliche Vorhersagekraft behalten. Die Auflösung liegt darin, dass sich nicht die Krümmung, sondern das Ziel verschiebt. Unter QAT hat sich das Netz bereits an den Quantisierungsfehler angepasst (\S\ref{sec:fund-quant}), sodass die isolierte Quantisierung einer einzelnen Schicht den Loss oft nicht mehr verschlechtert, sondern sogar leicht verbessert (\S\ref{sec:phase1-results}). Die Hessian-Krümmung bleibt also ein stabiles Merkmal der Netzarchitektur, ihre Interpretation als Schadens-Prädiktor setzt aber voraus, dass eine Störung tatsaechlich schadet, eine Voraussetzung, die PTQ erfüllt und QAT gerade durch sein Trainingsziel systematisch unterläuft. Praktisch bedeutet dies, dass eine Hessian-basierte Sensitivitätsanalyse für QAT-Modelle nicht ohne Weiteres von PTQ-Modellen übernommen werden sollte, selbst wenn die zugrunde liegende Krümmungslandschaft nahezu identisch aussieht.

\subsection{Deployment-Realität vs.\ theoretischer Vorteil}
\label{sec:disc-deployment}

Der in \S\ref{sec:fund-quant} beschriebene Bitshift-Vorteil von PoT ist eine Eigenschaft der Arithmetik und nicht automatisch eine Eigenschaft jeder verfügbaren Software-/Hardware-Kombination. Die Deployment-Ergebnisse (\S\ref{sec:phase10}) zeigen diesen Unterschied deutlich. Auf einer Standard-Rechenzentrums-GPU (A100) mit einem aktuellen, aber in seiner Faltungs-Kernel-Abdeckung eingeschränkten Quantisierungs-Framework (torchao 0.17) bringt die INT8-Quantisierung keinen Laufzeitvorteil und ist in der naheliegendsten Konfiguration deutlich langsamer als FP32. Auf einer CPU mit dedizierter INT8-Hardware-Unterstützung (AVX-512 VNNI) und einem Backend, das tatsächlich quantisierte Faltungs-Kernel implementiert (fbgemm), wird dagegen ein realer Durchsatzgewinn gemessen ($1.5$--$6.0\times$, CIFAR10, Tabelle~\ref{tab:fbgemmcpu}). Fbgemm beschleunigt generische Ganzzahl-Multiplikation, nicht den PoT-spezifischen Bitshift. Der theoretische Bitshift-Vorteil selbst bleibt in diesem Projekt auf keiner der beiden getesteten Plattformen gemessen, mangels eines nativen PoT-/Bitshift-Kernels.
\par\medskip

Die praktische Lehre daraus ist, dass es entscheidend von der Reife und dem Umfang der eingesetzten Inferenz-Software abhängt, ob ein theoretisch hardware-attraktives Quantisierungsschema tatsaechlich einen Vorteil liefert, und nicht nur vom Schema selbst. Die Speicherreduktion (theoretisch $4\times$ für Gewichte, \S\ref{sec:phase10}) bleibt der arithmetisch verlässlichere Vorteil, unabhängig davon, ob ein optimierter Rechenkernel existiert.

\subsection{Limitationen}
\label{sec:limitations}

\begin{itemize}
\item \textbf{Ein Trainings-Seed pro Kombination:} Für keine der sechs Architektur$\times$Datensatz-Kombinationen wurde über mehrere Trainings-Seeds gemittelt. Die Cross-Seed-Varianz der Top-Damage-Layer-Identifikation ist daher nicht charakterisiert; einzelne Schicht-Ränge könnten sich unter einem anderen Trainings-Seed verschieben. Als robust gilt das qualitative Muster, also die Top-Damage-Layer-Identität, der $S_{\mathrm{raw}}$-vs-$S_{\mathrm{hawq}}$-Vergleich und die top-$k^*$-Überlappungen, nicht die Einzelrang-Stabilität.

\item \textbf{Ungleiche Proben-Budgets zwischen Datensätzen:} CIFAR10 und ImageNet100 wurden mit unterschiedlichem Hutchinson-Proben-Budget vermessen, aus rechnerischer Notwendigkeit. Ein Teil der schwächeren Korrelation auf CIFAR10 (\S\ref{sec:phase4}) könnte auf Messrauschen statt auf einen echten Datensatzeffekt zurückgehen.

\item \textbf{Drei Architekturen, zwei Datensätze:} Die Generalisierbarkeit auf andere CNN-Familien (z.\,B.\ EfficientNet, MobileNet) oder auf nicht-Vision-Domänen ist ungetestet.

\item \textbf{Nur PoT:} Die Kernaussage, dass das HAWQ-V2-Produkt spezifisch durch PoT-Grid-Alignment-Effekte irregeführt wird (\S\ref{sec:disc-hawq}), ist eine PoT-spezifische Erklärung; ob und wie stark sie sich auf uniforme oder gemischt-präzise Schemata übertraegt, ist nicht getestet.

\item \textbf{PTQ als primärer Fokus:} QAT-Ergebnisse (\S\ref{sec:phase9}) werden berichtet, sind aber durchweg schwächer signifikant und methodisch komplizierter zu interpretieren, da isolierte QAT-Quantisierung den Loss teils verbessert statt verschlechtert.

\item \textbf{KFAC-Zerlegung gescheitert:} Die architektonische Ursachensuche über KFAC liefert keine belastbare Erklärung; die drei aufgestellten Hypothesen bleiben unentschieden.

\item \textbf{ImageNet100-Aufloesungskontrast für Phase~6 nicht auswertbar:} Der für eine Diskriminierung der KFAC-Hypothesen vorgesehene Kontrast zwischen den Eingabeauflösungen von CIFAR10 und ImageNet100 (\texttt{resolution\_contrast\_ratio}) ist in den verfügbaren Daten durchweg nicht bestimmt.

\item \textbf{fbgemm-CPU-Pfad nur auf CIFAR10 gemessen:} Der vollständige Durchsatz-/Genauigkeits-Sweep (\S\ref{sec:phase10}, Tabelle~\ref{tab:fbgemmcpu}) deckt alle drei Modelle und beide Quantisierungsstufen ab, aber nur auf CIFAR10. Für ImageNet100 liegt keine fbgemm-Messung vor. Zudem misst dieser Pfad INT8 im Allgemeinen, nicht PoTs Bitshift im Speziellen (\S\ref{sec:disc-deployment}); der PoT-spezifische Laufzeitvorteil bleibt auf beiden getesteten Plattformen ungemessen.
\end{itemize}

\subsection{Ausblick}
\label{sec:outlook}

Mehrere Anschlussrichtungen ergeben sich direkt aus den hier dokumentierten Negativresultaten und offenen Fragen. Erstens könnte eine RMT-informierte Zerlegung, die explizit zwischen Bulk- und Ausreißer-Eigenwerten des Hessian-Spektrums unterscheidet, dort ansetzen, wo die KFAC-Faktorisierung gescheitert ist, moeglicherweise mit einer gröberen, aber numerisch stabileren Zerlegung. Zweitens sollte die zentrale Erklärung dieses Projekts gegen uniforme und gemischt-präzise Quantisierungsschemata getestet werden, um zu prüfen, ob das HAWQ-V2-Scheitern tatsächlich PoT-spezifisch ist oder breiter auftritt. Drittens ist eine Ausweitung auf Transformer-Architekturen naheliegend, insbesondere da die einzige vorherige direkte Validierung der HAWQ-V2-Metrik \citep{hill2026saturationmakesquantizationerror} bereits im LLM-Setting stattfand (\S\ref{sec:related-work}); ein direkter CNN-vs-Transformer-Vergleich unter identischer Methodik wäre aufschlussreich. Viertens sollte eine Multi-Seed-Validierung die in \S\ref{sec:limitations} genannte Seed-Varianz-Lücke schließen.